% Complete additive module. Requires amsmath, amssymb, mathrsfs.
% Original coordinates, Fourier sign, Euler action and full arithmetic units
% are retained. No changes to the counterfactual owner's source files.
\section{All independent coefficient faces attached to the original theta source}
\label{sec:tau-mixed-coefficient-faces}

This construction attaches the independent coefficient masks of the original
mixed-support algebra to the actual theta complexes.  It retains the original
outer source label and leg mask independently of the coefficient mask.
A further slot-local leg diagram and its exact comparison with that outer
diagram are constructed explicitly.  The source-level maps,
their cohomology, and the full arithmetic boundary of the finite section are
computed below.  Every tensor operation already present in the ambient
programme keeps its original tensor variables; the finite direct sums in
this section are explicitly coefficient-slot observations.

\subsection{The original source and its full finite arithmetic sections}

Retain
\[
\begin{split}
V&=\{\phi\in\mathcal S(\mathbb R):\phi(-x)=\phi(x),\
                  \phi(0)=0,\ \int_{\mathbb R}\phi(x)\,dx=0\},\\
\mathscr B&=\{F\in C^\infty(\mathbb R_+):
 \sup_{x>0}x^b|D^jF(x)|<\infty\quad(b\in\mathbb Z,j\geq0)\},\\
D&=-x\partial_x,\qquad
\widehat\phi(\nu)=\int_{\mathbb R}\phi(x)e^{-2\pi ix\nu}\,dx,\\
\Theta\phi(x)&=2\sum_{n\geq1}\phi(nx),\qquad
Q=\mathscr B/\Theta V,\qquad q:\mathscr B\longrightarrow Q.
\end{split}
\tag{MCF1}
\]
The letter $q$ in (MCF1) denotes the quotient map; a coefficient Frobenius
prime power will be denoted $q_0$ below.  Put
\[
\begin{gathered}
U=\{s:0<\Re s<1,\ \Re s\ne\tfrac12\},\qquad
T_UQ=\bigcup_p\ker(p(D):Q\to Q),\\
\mathscr B_U=q^{-1}(T_UQ),\qquad
C_U=[V\xrightarrow\Theta\mathscr B_U],
\end{gathered}
\tag{MCF2}
\]
where $p$ ranges over all monic polynomials with roots in $U$, including $1$.
The original source calculation proves injectivity of $\Theta$,
$D\Theta=\Theta D$, and injectivity of every nonzero polynomial $p(D)$
on $V$ and $\mathscr B$.  We use its exact range identities
\[
\begin{gathered}
p(D)V=\ker(j_pH),\qquad
p(D)\mathscr B=\ker(j_p\mathcal M),\\
\mathcal MF(s)=\int_0^\infty F(x)x^s\frac{dx}{x},\qquad
\mathcal M\Theta\phi=gH_\phi,\qquad g=2\xi,\\
H_{P(D)\phi_*}=P,\qquad
\phi_*=(4\pi^2x^4-6\pi x^2)e^{-\pi x^2}.
\end{gathered}
\tag{MCF3}
\]
These are the original arithmetic identities (AG1) and (CAU1), with their
full-plane Euler inverses.  The attachment below uses their stated domains;
it does not introduce a different test-function completion.  For example,
the last identity follows from
\[
M_+\phi_*(s)=\tfrac12s(s-1)\pi^{-s/2}\Gamma(s/2),\qquad
H_\phi=\frac{2\pi^{s/2}M_+\phi}{s(s-1)\Gamma(s/2)}.
\]
The displayed Mellin formula is obtained by integrating the two Gaussian
monomials and using $\Gamma(z+1)=z\Gamma(z)$; the apparent exceptional
values are handled by the original entire continuation.  The moments of
$\phi_*$ are zero because $\phi_*(0)=0$ and
$4\pi^2\int x^4e^{-\pi x^2}dx-6\pi\int x^2e^{-\pi x^2}dx=3-3=0$.

For a nonempty finite full-order packet in $U$, write
\[
\begin{gathered}
h(s)=\prod_{\rho\in Z}(s-\rho)^{m_\rho},\qquad
m_\rho=\operatorname{ord}_\rho g,\qquad d_h=\deg h,\\
E_h=\mathbb C[s]/h,\quad A_h=M_s,\quad v_h=g/h,\quad
\upsilon_h=j_hv_h,\quad\varepsilon_h=\upsilon_h^{-1},\\
F_h=S_h^{\mathscr B}\Theta\phi_*,\qquad
P_h(u)=r_h(\varepsilon_hu),\qquad
s_hu=P_h(u)(D)F_h,\quad \sigma_h=qs_h.
\end{gathered}
\tag{MCF4}
\]
Here $r_h$ is the unique polynomial representative of degree less than
$d_h$, and $S_h^{\mathscr B}$ is the inverse of $h(D)$ on the exact range
in (MCF3).  Every actual multiplicity is retained, so $v_h$ is nonzero
at each packet centre and $\upsilon_h$ is a unit in the full local algebra.
Every packet index and directed limit in this module also includes the
empty packet $h=1$, with
$F_1=\Theta\phi_*$, $E_1=0$, $s_1=\sigma_1=\ell_1=0$, and
$\mathscr B_1=\Theta V$.  Its maps into other packet stages are zero
on $E_1$ and the original inclusion on $\Theta V$.  No inverse unit
at this zero-vector-space stage is used.  The definition of $F_1$ is
forced by $1(D)F_1=\Theta\phi_*$ and gives $\mathcal MF_1=g$.
Thus the directed systems
are defined even when there is no nonempty packet in $U$.
The equalities
\[
\begin{gathered}
h(D)F_h=\Theta\phi_*,\qquad \mathcal MF_h=v_h,\qquad
J_hs_h=1,\quad J_h=j_h\mathcal M,\qquad J_h\Theta=0,\\
h(D)s_hu=\Theta(P_h(u)(D)\phi_*),\\
Ds_h-s_hA_h=\Theta\phi_*\ell_h,\qquad
\ell_h(u)=[s^{d_h-1}]P_h(u)
\end{gathered}
\tag{MCF5}
\]
have the original source types.  To verify the last equality in full,
$sP_h(u)-P_h(A_hu)$ has zero $h$-jet, degree at most $d_h$, and leading
coefficient $\ell_h(u)$.  Monicity gives
$sP_h(u)-P_h(A_hu)=h\ell_h(u)$.  Apply this polynomial in $D$ to $F_h$
and use the first equality.  The other identities follow directly from
(MCF3)--(MCF4) and retain the entire multiplier $v_h$.

For completeness the finite arithmetic image is exactly
$\sigma_hE_h=Q[h(D)]$.  If $h(D)F=\Theta\phi$, the assignment
$[F]\mapsto j_hH_\phi$ is well-defined because replacing $F$ by
$F+\Theta\psi$ replaces $\phi$ by $\phi+h(D)\psi$.  A zero jet makes
$\phi=h(D)\psi$ by (MCF3), and injectivity of $h(D)$ gives
$F=\Theta\psi$.  Every jet represented by a polynomial $P$ is obtained
from $S_h^{\mathscr B}\Theta(P(D)\phi_*)$, since $j_h(gP)=0$.
Thus this assignment is an isomorphism.  Applied to $s_hu$ it gives
$\varepsilon_hu$ by (MCF5), an invertible change of full jet coordinates.
In particular $\sigma_h$ is injective and has the asserted image.

\subsection{The independent coefficient face category and its exact quotient}

Fix $n\geq1$ and $I_n=\{0,\ldots,n-1\}$.  For every subset
$A\subseteq I_n$, including the empty subset, put
\[
C_U[A]=\bigoplus_{i\in A}C_U
 =[V^A\xrightarrow{(\phi_i)\mapsto(\Theta\phi_i)}\mathscr B_U^A],
\qquad E_h[A]=E_h^A.
\tag{MCF6}
\]
Empty direct sums are zero vector spaces.  If $A\subseteq B$, define
$z_{AB}$ by the identity in each old coordinate and by the typed zero in
every coordinate of $B\setminus A$, in both degrees and in $E_h[A]$.
Then $d z_{AB}=z_{AB}d$ coordinate by coordinate, and
$z_{BC}z_{AB}=z_{AC}$.  This is a functor from the inclusion category
$\mathcal P(I_n)$ to complexes.  The complete face diagram is retained;
no colimit is taken that would identify its different masks.

For a complex vector space $X$, set
$G(X)=\{\tau\}\sqcup\{x^\bullet:x\in X\}$, with the original supported
addition and scalar action.  The degree-$j$ mixed carrier is
\[
\mathscr D_n^j=\prod_{i\in I_n}G(C_U^j)
\simeq\coprod_{A\subseteq I_n}\{A\}\times(C_U^j)^A.
\tag{MCF7}
\]
The empty term on the right is one point, corresponding to the all-absent
tuple.  The bijection sends each present coordinate to its original
amplitude and records every present zero in $A$.  Its inverse puts those
amplitudes in the coordinates of $A$ as supported values and puts $\tau$
in the other coordinates.  Thus all $2^n$ masks, including the proper
nonempty faces, are present.  A face arrow $z_{AB}$ changes the recorded
mask from $A$ to $B$ and fills the new coordinates with supported zeros.
It is an arrow between labelled faces, not an identification of them.

The differential preserves the original mask of its input and applies
$\Theta$ independently in every present coordinate.  A source vector
annihilated by its amplitude differential reaches the supported zero at
the same mask.  In the internal split cycle and boundary construction,
$\epsilon(A,x)=(A,0)$ and
$Z=\operatorname{Eq}(d,\epsilon d)$; boundary relations are the
coequalizer of $b$ and $\epsilon b$ within each labelled fibre.
Consequently the exact cohomology and quotient are
\[
\begin{gathered}
H^0(C_U[A])=0,\qquad H^1(C_U[A])=(T_UQ)^A,\\
\mathscr D_n^1\longrightarrow\prod_iG(T_UQ),\qquad
F_i^\bullet\longmapsto(qF_i)^\bullet,\quad\tau\longmapsto\tau.
\end{gathered}
\tag{MCF8}
\]
Indeed $\Theta$ is injective in every active coordinate, and the image
in degree one is exactly $(\Theta V)^A$.  Two amplitudes at mask $A$
are identified precisely when their coordinate differences belong to
$(\Theta V)^A$.  The fibre over the supported zero at $A$ is
$\{A\}\times(\Theta V)^A$; it is not combined with any other mask.
More generally, choosing any lift $F$ of an amplitude in $(T_UQ)^A$
gives its entire fibre as $F+(\Theta V)^A$ at that same mask.

The full finite arithmetic attachment is
\[
\begin{array}{ccccc}
E_h[A]&\xrightarrow{\ s_h^A\ }&\mathscr B_U^A
       &\xrightarrow{\ q^A\ }&(T_UQ)^A\\
z_{AB}\downarrow&&z_{AB}\downarrow&&\downarrow z_{AB}\\
E_h[B]&\xrightarrow{\ s_h^B\ }&\mathscr B_U^B
       &\xrightarrow{\ q^B\ }&(T_UQ)^B.
\end{array}
\tag{MCF9}
\]
All squares commute because every displayed map is applied independently
to the unchanged old coordinates and is linear on the inserted zeros.
Every $\sigma_h^A=q^As_h^A$ is injective by (MCF5).  In every active
coordinate its exact Euler boundary is
\[
(D^As_h^A-s_h^AA_h^A)(u_i)_{i\in A}
 =(\Theta\phi_*\ell_h(u_i))_{i\in A}.
\tag{MCF10}
\]
This is an equality of actual functions, whose primitive is
$(\phi_*\ell_h(u_i))_{i\in A}\in V^A$.
After the original quotient it proves $D^A\sigma_h^A=\sigma_h^AA_h^A$.
Thus every proper coefficient face carries the same full arithmetic
boundary calculation, with its mask still specified.

\subsection{Independent slot-local leg presentations}

There is also a slot-local copy of the original two-leg presentation at
every active coefficient.  This auxiliary indexing does not replace the
original outer leg label: its exact attachment, including empty coefficient
support, is (MCF34)--(MCF40) below.  Its objects are
\[
\lambda=(A,(S_i)_{i\in A}),\qquad
A\subseteq I_n,\quad\varnothing\ne S_i\subseteq\{+,-\}.
\tag{MCF11}
\]
For empty $A$ the family is empty.  Order these objects by
$\lambda\leq\lambda'$ when $A\subseteq A'$ and $S_i\subseteq S'_i$
for every $i\in A$.  Thus each coefficient slot has four states: absent,
plus only, minus only, or both.  There are exactly $4^n$ slot-local masks.

At $\lambda$ put
\[
\begin{split}
C_\lambda^0&=\bigoplus_{i\in A}
   \left(\bigoplus_{+\in S_i}V_+\ \oplus\!
                         \bigoplus_{-\in S_i}V_-\right),\\
C_\lambda^1&=\mathscr B_U^A,\qquad
d_\lambda(\phi_i,\psi_i)_i
       =(\Theta\phi_i-\Theta\widehat\psi_i)_i.
\end{split}
\tag{MCF12}
\]
A coordinate of a missing leg in the displayed formula is its typed zero;
that convention evaluates the formula without adjoining that leg to the
mask.  The plus and target actions of $t\in\mathbb C[t]$ are $D$; the
minus action is $1-D$.  The Fourier identity
$\widehat{D\phi}=(1-D)\widehat\phi$ proves that this is a complex of
$\mathbb C[t]$-modules.  For $\lambda\leq\lambda'$, retain every old
source leg and coefficient, insert typed zeros in the new legs and slots,
and use zero insertion in degree one.  Both differentials then agree
term by term, proving a strictly composing cochain map.  The actual
supported arrow changes the label from $\lambda$ to $\lambda'$ and
retains every newly supported zero as such.

The natural comparison with (MCF6) is the cochain map
\[
p_\lambda:C_\lambda\longrightarrow C_U[A],\qquad
p_\lambda^0(\phi_i,\psi_i)_i
       =(\phi_i-\widehat\psi_i)_i,\qquad
p_\lambda^1=1.
\tag{MCF13}
\]
It commutes with every zero-insertion transition.  In a plus-only slot
its degree-zero map is the identity, in a minus-only slot it is
$-\widehat{\phantom\phi}$, and in a joint slot it has the exact inverse
coordinates
\[
u=\phi-\widehat\psi,\quad v=\tfrac12(\phi+\widehat\psi),\qquad
\phi=v+\tfrac12u,\quad\psi=\widehat v-\tfrac12\widehat u.
\tag{MCF14}
\]
Fourier squares to the identity on $V$, so every displayed map and
inverse has its stated domain.  Write
$J(\lambda)=\{i\in A:S_i=\{+,-\}\}$.  The exact kernel and cohomology are
\[
\begin{gathered}
0\longrightarrow V^{J(\lambda)}[0]
\xrightarrow{\delta_\lambda} C_\lambda
\xrightarrow{p_\lambda} C_U[A]\longrightarrow0,\\
\delta(v)_i=(v_i,\widehat v_i)\quad(i\in J(\lambda)),\\
H^0(C_\lambda)=V^{J(\lambda)},\qquad
H^1(C_\lambda)=(T_UQ)^A.
\end{gathered}
\tag{MCF15}
\]
Surjectivity of $p^0$ follows from the single-leg formulas and (MCF14).
Its kernel in a joint slot is exactly the displayed diagonal, and its
degree-one kernel is zero.  Its differential on $(u,v)$ is $\Theta u$.
Injectivity of $\Theta$ now proves the cohomology statement directly.
Thus the two-leg diagonal remains present at every relevant independent
coefficient; it is not removed when the arithmetic degree-one quotient
is observed.

The slot-local cohomology carrier is the following.  The separate outer
label is retained in (MCF34).
\[
\widetilde H_n^1
=\{\tau_n\}\sqcup
 \coprod_{\substack{\lambda=(A,S)\\ A\ne\varnothing}}
               \{\lambda\}\times(T_UQ)^A.
\tag{MCF16}
\]
Forgetting only the leg masks gives a pointed map to $\prod_iG(T_UQ)$.
Over a specified tuple of coefficient mask $A$ it has exactly
$3^{|A|}$ labelled preimages, one for every nonempty $S_i$; the amplitudes
are fixed.  Over the all-absent tuple its fibre is the single $\tau_n$.
A specified section takes the joint leg mask at every active coefficient.
These formulas prove the exact relationship between the combined support
diagram and the coefficient-mask observation; neither is silently used
in place of the other.  The maps $s_h^A$ and $\sigma_h^A$ of (MCF9)
give their source and arithmetic amplitudes at every $\lambda$, with the
same recorded coefficient and leg masks.

One can also lift each original one-leg complex into its selected leg
presentation.  At a single slot the degree-zero section of $p_\lambda$ is
\[
j_+(u)=u,\qquad j_-(u)=-\widehat u,\qquad
j_{+-}(u)=(u/2,-\widehat u/2),
\tag{MCF17}
\]
and degree one is the identity.  They are $\mathbb C[t]$-linear cochain
maps by the Fourier identity.  For the two inclusions into the joint
slot their precise difference is
\[
\operatorname{inc}_+j_+-j_{+-}=\delta(u/2),\qquad
\operatorname{inc}_-j_--j_{+-}=-\delta(u/2)
\quad\hbox{in degree zero},
\tag{MCF18}
\]
and zero in degree one.  These retain the original diagonal; a claim of
strict naturality for these three chosen sections would omit it.
There is also an explicit complex-linear homotopy.  The coherent original
sections give $\mathscr B_U=\Theta V\oplus s_UE_U$, with
$E_U=\varinjlim_hE_h$; define $t_U(\Theta\phi+s_Uu)=\phi$.
Then $t_U\Theta=1$.  The degree-minus-one maps
$H_\pm(F)=\pm\delta(t_UF)/2$ obey $dH_\pm=0$ and
$H_\pm\Theta=\pm\delta/2$, so $dH_\pm+H_\pm d$ equals (MCF18).
This is the exact homotopy of the chosen sections.  No equivariance of
$t_U$ or $H_\pm$ is asserted.  Existence of the coherent $s_U$ and its
original unit is checked in the divisor calculation below.

\subsection{Arithmetic dilation on every face and its full source boundary}

For $a>0$ write $\mathcal R_aF(x)=F(x/a)$, with the same formula on $V$.
Changing variables shows that it preserves the original spaces, preserves
the value-zero condition, and multiplies the source integral by $a$.
Direct substitution gives
\[
\Theta\mathcal R_a=\mathcal R_a\Theta,\qquad
\widehat{\mathcal R_a\phi}=a\mathcal R_{1/a}\widehat\phi,\qquad
\mathcal M\mathcal R_aF=a^s\mathcal MF.
\tag{MCF19}
\]
It commutes with $D$ and hence preserves $T_UQ$ and its inverse image
$\mathscr B_U$.  On every combined face its action is
\[
(\phi_i,\psi_i)_i\longmapsto
 (\mathcal R_a\phi_i,a\mathcal R_{1/a}\psi_i)_i,\qquad
F_i\longmapsto\mathcal R_aF_i.
\tag{MCF20}
\]
The coefficient and leg masks are unchanged.  The Fourier formula in
(MCF19) makes the two differentials commute exactly.  These actions
also commute with every face inclusion and with $p_\lambda$, since
\[
\mathcal R_a\phi-\widehat{a\mathcal R_{1/a}\psi}
      =\mathcal R_a(\phi-\widehat\psi).
\]
On the diagonal the action is $v\mapsto\mathcal R_av$, because its
minus coordinate becomes $a\mathcal R_{1/a}\widehat v
=\widehat{\mathcal R_av}$.  Thus (MCF15) is an equivariant exact sequence.

The finite arithmetic action is
\[
A_{a,h}=M_{j_h(a^s)}:E_h\longrightarrow E_h,\qquad
A_{a,h}|_\rho=a^\rho\sum_{j=0}^{m_\rho-1}
             \frac{(\log a)^j}{j!}N_\rho^j,\quad N_\rho=M_{s-\rho}.
\tag{MCF21}
\]
The complete nilpotent Taylor polynomial is displayed.  Mellin and (MCF19)
give $J_h\mathcal R_a=A_{a,h}J_h$.  Dilation preserves $Q[h(D)]$, and
$J_h$ is the inverse of $\sigma_h$ there by (MCF5).  Therefore
$\overline{\mathcal R}_a\sigma_h=\sigma_hA_{a,h}$, an exact equality
on the original arithmetic quotient.

The source defect has a full original primitive, rather than an assumed
equivariant lift.  Define
\[
\begin{split}
\chi_{a,h}(u)&=\mathcal R_a(P_h(u)(D)\phi_*)
                       -P_h(A_{a,h}u)(D)\phi_*\in V,\\
\Psi_{a,h}(u)&=S_h^V\chi_{a,h}(u),
\end{split}
\tag{MCF22}
\]
where $S_h^V$ is the inverse of $h(D)$ on $\ker(j_hH)$ from (MCF3).
This inverse is applied on its proved domain: the Mellin calculation gives
\[
H_{\chi_{a,h}(u)}=a^sP_h(u)-P_h(A_{a,h}u),\qquad
j_hH_{\chi_{a,h}(u)}
 =j_h(a^s)\varepsilon_hu-\varepsilon_hA_{a,h}u=0.
\tag{MCF23}
\]
The complete $\varepsilon_h=j_h(g/h)^{-1}$ occurs in both terms.
Applying $h(D)$ to the two sides and using (MCF5), (MCF19), and
$h(D)\Psi=\chi$, gives
\[
\mathcal R_as_hu-s_hA_{a,h}u=\Theta\Psi_{a,h}(u).
\tag{MCF24}
\]
Indeed the difference after application of $h(D)$ is zero, and its
injectivity on $\mathscr B$ proves the displayed equality.  In every
coefficient face this equality is applied separately to each $u_i$.
In a combined leg presentation an actual degree-zero primitive is
$j_{S_i}(\Psi_{a,h}(u_i))$ from (MCF17).  Its differential is precisely
the corresponding coordinate of (MCF24).  The alternative single-leg
and joint-leg primitives differ by the diagonal terms of (MCF18),
which remain recorded in (MCF15).

The primitives themselves have the exact composition law
\[
\Psi_{ab,h}(u)=\mathcal R_a\Psi_{b,h}(u)
                         +\Psi_{a,h}(A_{b,h}u).
\tag{MCF25}
\]
Expand $\mathcal R_{ab}s_h-s_hA_{a,h}A_{b,h}$ as
$\mathcal R_a(\mathcal R_bs_h-s_hA_{b,h})
 +(\mathcal R_as_h-s_hA_{a,h})A_{b,h}$, apply (MCF24), and cancel the
injective $\Theta$.  This proves the law with the original dilation
order, source action, and full arithmetic multiplier.

\subsection{Packet enlargement with every coefficient face retained}

For full-order packets $h\mid H$, put $b=H/h$.  Their supports differ by
whole local factors, so $b$ and $h$ are coprime.  Let
$\iota_{hH}:E_h\to E_H$ insert the old full jets and zero jets at every
added centre.  CRT gives this injective $\mathbb C[s]$-linear map and
the restriction left inverse.  At old centres $\upsilon_H=b^{-1}\upsilon_h$
and $\varepsilon_H=b\varepsilon_h$.  At added centres $b$ vanishes to
the full required order.  Hence
\[
r_H(\varepsilon_H\iota_{hH}u)=b\,r_h(\varepsilon_hu),\qquad
b(D)F_H=F_h,\qquad s_H\iota_{hH}=s_h.
\tag{MCF26}
\]
When $h=1$, $b=H$ and the middle identity is exactly
$H(D)F_H=\Theta\phi_*=F_1$; the first and third identities have
zero domain and are zero-map identities.  The highest-coefficient,
action, and primitive identities below also have zero domain there.
For a nonempty $h$, the first identity holds in every full local factor and both polynomials
have degree less than $\deg H$; uniqueness of the remainder proves it.
For the second identity, apply $h(D)$ and use the injectivity in (MCF3).
Together they prove the third.  Monicity of $b$ preserves the highest
coefficient of the remainders in the first identity, giving
\[
\ell_H\iota_{hH}=\ell_h,\qquad
A_{a,H}\iota_{hH}=\iota_{hH}A_{a,h},\qquad
\Psi_{a,H}\iota_{hH}=\Psi_{a,h}.
\tag{MCF27}
\]
The action identity follows in the unchanged old jets and the zero added
jets.  The primitive identity follows from (MCF24), the section identity,
and injectivity of $\Theta$.  All statements hold coordinatewise in
$E_h[A]$ and commute with every coefficient and combined-leg face map.
Thus the finite face attachments glue by literal source equalities.
Their directed union supplies the coherent $s_U$ used in (MCF18).
The decomposition $\mathscr B_U=\Theta V\oplus s_UE_U$ follows from
surjectivity onto $T_UQ$ and injectivity of each $\sigma_h$; every vector
and every finite relation occurs in a common full packet.  The original
all-polynomial calculation (CAU5)--(CAU13) proves that these packets
exhaust $T_UQ$, retaining every local zero order.

\subsection{The original negacyclic masks and coefficient Frobenius}

At its finite arithmetic-algebra vertex, the complete independent carrier
is the original
\[
D_{E_h,n}=G(E_h)^n,\qquad B_{h,n}=E_h[t]/(t^n+1),\qquad
\pi_n(c)=\sum_{i=0}^{n-1}p(c_i)t^i.
\tag{MCF28}
\]
The variable $t$ here is the additional coefficient variable of this
finite observation.  The spectral variable in $E_h$ is still $s$.
The negacyclic product, its original unit, and its mask law are
\[
\begin{gathered}
(c\star d)_k=\mathop{\bigoplus}_{i+j=k}c_id_j\ \oplus\!
             \mathop{\bigoplus}_{i+j=k+n}(-1)^\bullet c_id_j,
\qquad 1_D=(1^\bullet,\tau,\ldots,\tau),\\
\operatorname{mask}(c+d)=\operatorname{mask}(c)\cup\operatorname{mask}(d),\qquad
\operatorname{mask}(c\star d)
 =\operatorname{mask}(c)+\operatorname{mask}(d)\quad(\bmod n).
\end{gathered}
\tag{MCF29}
\]
For the product, an output is supported exactly when at least one pair
of supported input indices sums to it.  Wrap signs are supported and a
nonempty sum of supported summands remains supported even when its
coefficient cancels.  This proves the stated law with every mixed face.
Monic remainder in $t$ gives the unique coefficients of $\pi_n(c)$;
the mask then distinguishes each coefficient zero from absence.  Thus
$(\operatorname{mask},\pi_n)$ is injective.  Its facewise source and
cohomology attachment is exactly (MCF9), or (MCF16) when leg masks are
also retained.  The product in (MCF29) belongs to this finite algebra
vertex.  The source maps $s_h^A$ have their stated linear types; no
unproved ring multiplication on the original source is inserted.

Let $q_0$ be an odd prime power with $\gcd(q_0,n)=1$, and put
$\kappa(i)=q_0i\bmod n$ and
$\epsilon_i=(-1)^{\lfloor q_0i/n\rfloor}$.  The coefficient Frobenius
lift on the entire independent object is
\[
c_i\longmapsto\epsilon_i^\bullet c_i
       \quad\hbox{in slot }\kappa(i),\qquad
A\longmapsto\kappa(A).
\tag{MCF30}
\]
An absent coordinate stays absent.  On a combined label it sends
$S_i$ to $S'_{\kappa(i)}=S_i$ and multiplies every present source-leg
and target amplitude in that coefficient by $\epsilon_i$.  It commutes
with the differential because it uses the same coefficient sign in
each term of that differential.  It commutes with face inclusion,
the maps $p_\lambda$, $s_h^A$, $\sigma_h^A$, and all arithmetic dilations
by coordinate permutation and complex linearity.

On the integral coefficient ring $\mathbb Z[t]/(t^n+1)$ this is the
matrix of $t\mapsto t^{q_0}$.  It is well-defined since
$(t^{q_0})^n=-1$, and choosing $r$ with $q_0^r\equiv1\pmod{2n}$
proves its $r$th power is the identity.  Reduction to $\mathbb F_{q_0}$
is exactly coefficient Frobenius; extension to $\mathbb C$ retains the
same integral signed-permutation matrix.  On supported monomials,
reducing the exponent before or after substitution gives the same
wrap sign, since $q_0$ is odd.  Distributivity and the absent cases
therefore prove multiplicativity on (MCF29), before synchronization.
In particular its support action is the exact permutation in (MCF30),
rather than an observation of one global Boolean support bit.

The synchronization arrow is also retained explicitly:
\[
\mathsf E_n=(1^\bullet,e,\ldots,e),\qquad
\mathsf E_nD_{E_h,n}\simeq G(B_{h,n}),\qquad
D_{E_h,n}[\mathsf E_n^{-1}]\simeq\mathsf E_nD_{E_h,n}.
\tag{MCF31}
\]
Its multiplication fills every missing coordinate of an active input
with a supported zero while leaving all ordinary coefficients unchanged;
every output has a supported summand using any chosen active input
coordinate.  Hence $\mathsf E_n^2=\mathsf E_n$.  A unital map making
this idempotent invertible sends it to 1 and identifies $c$ with
$\mathsf E_nc$, proving the localization universal property.
If $f=\sum f_it^i$ and $J(f)=\{i:f_i\ne0\}$, its fibre over
$f^\bullet$ consists of one original vector for each nonempty mask
$A$ with $J(f)\subseteq A\subseteq I_n$.  The fibre over $\tau$ is
the all-absent vector alone.  In particular a supported zero has
$2^n-1$ distinct independent-mask preimages.  These exact fibres remain
attached; the synchronized observation does not replace (MCF28).

Finally the original arithmetic coefficients have an exact retract
\[
\iota_n:E_h\longrightarrow B_{h,n},\quad f\mapsto f,
\qquad r_n:B_{h,n}\longrightarrow E_h,
\quad r_n(b)=\frac1n\operatorname{Tr}_{B_{h,n}/E_h}(M_b),
\qquad r_n\iota_n=1.
\tag{MCF32}
\]
In the basis $1,t,\ldots,t^{n-1}$, multiplication by $t^i$ has no
diagonal entries for $1\leq i<n$, while multiplication by $b_0$
has $n$ diagonal entries $b_0$.  Thus $r_n(\sum b_it^i)=b_0$ and
its kernel is exactly $\bigoplus_{i=1}^{n-1}E_ht^i$.
Writing $\Psi_{q_0}$ for (MCF30) on coefficients, it fixes $\iota_n$
and preserves $r_n$.  The coefficientwise extension of $A_{a,h}$
commutes with it and intertwines both arrows in (MCF32).  On a coefficient
eigenspace of eigenvalue $\omega$, define the combined operator
$T_{a,q_0}=A_{a,h}\Psi_{q_0}=\Psi_{q_0}A_{a,h}$ on $B_{h,n}$.
Its actual full local action is
\[
\omega a^\rho\sum_{j=0}^{m_\rho-1}
          \frac{(\log a)^j}{j!}N_\rho^j.
\tag{MCF33}
\]
Every $\omega$ is a root of unity by the proved finite order, and the
original arithmetic action is recovered on the constant-coefficient
retract.  This completes the attachment of independent mixed coefficient
support to the original theta source, its combined leg diagram, and its
actual arithmetic operators.  All masks, source diagonals, full units,
nilpotents, and comparison fibres have their displayed maps.

\subsection{The independent outer source label, including empty coefficient support}

The original outer source label is $(W,S)$, with $W\subseteq V$ a
complex linear $D$-invariant subspace and
$\varnothing\ne S\subseteq\{+,-\}$.  It is independent
of the coefficient mask $A\subseteq I_n$.  Denote the set of those subspaces
$W$ by $\mathcal W_D$.  Its exact coefficient diagram is
\[
\begin{split}
\mathcal J_n^{\rm out}&=\{(W,S;A): W\in\mathcal W_D,
       \ \varnothing\ne S\subseteq\{+,-\},\ A\subseteq I_n\},\\
C_{W,S;A}^0&=
 \left(\bigoplus_{+\in S}W\ \oplus\!
                          \bigoplus_{-\in S}\widehat W\right)^A,
\qquad C_{W,S;A}^1=\mathscr B_U^A,\\
d_{W,S;A}(\phi_i,\psi_i)_i
      &=(\Theta\phi_i-\Theta\widehat\psi_i)_i.
\end{split}
\tag{MCF34}
\]
The minus leg is exactly $\widehat W$; it is not replaced by $W$ unless
the two spaces are equal.  Its action is $1-D$, preserving that space
by the Fourier identity.  The target restriction $\mathscr B_U\subseteq
\mathscr B$ is an injective cochain map to the original outer complex,
since $\Theta W\subseteq\Theta V\subseteq\mathscr B_U$.
For $W\subseteq W'$, $S\subseteq S'$ and $A\subseteq A'$, the
transition includes the old source subspaces and coordinates and inserts
the typed zeros in new legs and coefficient slots.  The degree-one
transition is zero insertion.  These are cochain maps, by substitution
in the displayed differential, and compose literally.

At $A=\varnothing$ this is a zero complex \emph{at its retained outer
label $(W,S)$}.  Its single coefficient value has empty coefficient
support; the outer leg mask is still $S$.  In the represented carrier
one retains all triples $(W,S;A,x)$ and separately adjoins a global
external point $\tau_{\rm out}$.  In particular
\[
(W,S;\varnothing,0)\ne\tau_{\rm out},\qquad
(W,S;A,0)\ne(W,S;A',0)\quad\text{if }A\ne A'.
\tag{MCF35}
\]
These are explicit distinct elements of the disjoint labelled carrier;
they do not assert a nonzero amplitude in a zero complex.  Taking a
colimit over coefficient masks would identify these zero insertions
with their images in the terminal full face.  That colimit is not part
of this construction.

Define $Q_{W,U}=\mathscr B_U/\Theta W$ and
$\eta_W:V/W\to Q_{W,U}$ by $[v]\mapsto[\Theta v]_W$.
The complete arithmetic quotient and kernel are
\[
\begin{gathered}
0\longrightarrow(V/W)^A\xrightarrow{\eta_W^A}Q_{W,U}^A
       \xrightarrow{\pi_W^A}(T_UQ)^A\longrightarrow0,\\
H^0(C_{W,S;A})=
 \begin{cases}W^A,&S=\{+,-\},\\0,&|S|=1,\end{cases}
\qquad H^1(C_{W,S;A})=Q_{W,U}^A.
\end{gathered}
\tag{MCF36}
\]
For the kernel sequence, injectivity of $\Theta$ proves that
$\Theta v\in\Theta W$ exactly when $v\in W$.  A class of
$\mathscr B_U/\Theta W$ maps to zero modulo $\Theta V$ exactly when
it is represented by $\Theta v$ for some $v\in V$.  Surjectivity
follows from the definition of $\mathscr B_U$.  Applying this proof
independently in the coordinates proves the exact sequence.
For the cohomology calculation, a joint slot has the reversible
coordinates (MCF14) with $u,v\in W$, since $\psi\in\widehat W$.
Its differential is $\Theta u$ and its diagonal is $W$.
A single-leg differential is injective with image $\Theta W$.
This proves every case, including the zero coefficient vector spaces
at empty $A$.  For $W\subseteq W'$ the transition on degree one has
kernel exactly $(W'/W)^A$ under the same original map $[v]\mapsto
[\Theta v]_W$.  The outer-label kernel is thus fully retained.

The exact relation to the slot-local diagram is an enlarged index
\[
(W,S;A,(T_i)_{i\in A}),\qquad
\varnothing\ne T_i\subseteq S.
\tag{MCF37}
\]
Its source uses $W$ on the present plus legs and $\widehat W$ on the
present minus legs, according to $T_i$, and its target is
$\mathscr B_U^A$.  Every differential is the original signed theta
differential in that coordinate.  The outer diagram (MCF34) embeds
by $T_i=S$ for every $i\in A$, with the identity on its coefficient
spaces.  This embedding remains injective on the index when $A$ is
empty because $W$ and $S$ are retained separately.
Forgetting $W,S$ and including $W,\widehat W$ into $V$ gives the
explicit cochain comparison to (MCF12); its target map is the identity
on $\mathscr B_U^A$.  Its induced map on $H^1$ is precisely $\pi_W^A$ in
(MCF36), and its induced map on $H^0$ is the diagonal inclusion
$W^{\{i:T_i=\{+,-\}\}}\subseteq V^{\{i:T_i=\{+,-\}\}}$.
For fixed $W=V$, forgetting only the outer $S$ has exactly three
labelled preimages over the empty slot-local mask, two when the union
of the nonempty $T_i$ is a single leg, and one when that union contains
both legs.  Indeed these are exactly the nonempty subsets
$S\subseteq\{+,-\}$ containing $\bigcup_iT_i$.
This computes the lost outer information of that specified forgetting
map; its empty coefficient value is not identified with global
absence in (MCF35).

There is a useful original-outer-source retract, which also checks the
empty coefficient endpoint without changing the label.  At a fixed
$(W,S)$, write $C_{W,S}^U$ for its one-coefficient original complex
with central term $\mathscr B_U$, and take the constant coefficient map
\[
c_0:C_{W,S;A}\longrightarrow C_{W,S}^U,\qquad
c_0(x)=\begin{cases}x_0,&0\in A,\\0,&0\notin A,
\end{cases}
\qquad
i_0:C_{W,S}^U\longrightarrow C_{W,S;\{0\}},\quad x\longmapsto(x)_0.
\tag{MCF38}
\]
Here $i_0$ includes the original outer complex as the coefficient-face
vertex $A=\{0\}$ in the full represented diagram.  This notation does
not assert an inclusion into every other face.  Each $c_0$ is a cochain
map and commutes with all coefficient zero insertions, since a newly
inserted coordinate has zero amplitude.  On the represented carrier
it sends $(W,S;\varnothing,0)$ to the supported outer zero at
$(W,S)$.  On global $\tau_{\rm out}$ the pointed extension
has value $\tau_{\rm out}$.  Its composite with the selected $i_0$
is the identity on the original outer coefficient amplitude and label.
The coefficient Frobenius fixes index zero and its sign is one, so
$c_0\Psi_{q_0}=c_0$.  This proves an actual outer-label-preserving
comparison, rather than identifying the original outer label with a
slot-local mask.

The finite arithmetic section at a proper outer source is
$\sigma_{h,W}(u)=[s_hu]_W\in Q_{W,U}$.  It is injective because
$\pi_W\sigma_{h,W}=\sigma_h$ and $\sigma_h$ is injective.
Its exact retained defects are
\[
\begin{gathered}
D\sigma_{h,W}-\sigma_{h,W}A_h
       =\eta_W([\phi_*\ell_h(\,\cdot\,)]_W),\\
\overline{\mathcal R}_a^{\,W\to\mathcal R_aW}\sigma_{h,W}
       -\sigma_{h,\mathcal R_aW}A_{a,h}
       =\eta_{\mathcal R_aW}([\Psi_{a,h}(\,\cdot\,)]_{\mathcal R_aW}).
\end{gathered}
\tag{MCF39}
\]
The first is (MCF5) in the quotient by $\Theta W$ and the second is
(MCF24) in the quotient by $\Theta\mathcal R_aW$.  Dilation transports
the outer label to $\mathcal R_aW$, as required by (MCF19); it need
not fix an arbitrary $D$-invariant $W$.  On the minus source it maps
$\widehat W$ to $a\mathcal R_{1/a}\widehat W=
\widehat{\mathcal R_aW}$.  This proves that every arrow in (MCF39)
has the stated domain and codomain.  The original primitives belong
to $V$; their possibly nonzero classes in $V/W$ or
$V/\mathcal R_aW$ are retained in (MCF36).  No primitive in $V$ has
been declared to be a boundary in a smaller source $W$.
For the full outer source $W=V$, both kernel classes are zero and
(MCF39) recovers the proved full-source intertwining.

Apply (MCF39) coordinatewise on every $A$, keeping all outer and
slot-local labels.  Packet inclusion, zero insertion, and coefficient
Frobenius commute with these defects by (MCF27), linearity and the
permutation (MCF30).  In particular the complete retained data form
the commuting represented diagram
\[
\{C_{W,S;A}\}_{W,S,A} \longrightarrow\
\{Q_{W,U}^A[-1]\}_{W,S,A}\ \xrightarrow{\{\pi_W^A\}}\
\{(T_UQ)^A[-1]\}_{W,S,A},\qquad
\ker\pi_W^A=\eta_W^A((V/W)^A),
\tag{MCF40}
\]
with the enriched slot-local faces and their inclusion and forgetting
maps specified in (MCF37).  The first arrow is the actual cochain map
to $Q_{W,U}^A[-1]$ given by zero in degree zero and the original
quotient in degree one.  Its exact kernel complex is also retained:
\[
K_{W,S;A}=[C_{W,S;A}^0\xrightarrow{d}(\Theta W)^A],\qquad
0\longrightarrow K_{W,S;A}\longrightarrow C_{W,S;A}
\longrightarrow Q_{W,U}^A[-1]\longrightarrow0.
\tag{MCF40a}
\]
Degreewise kernel computation proves this short exact sequence.
The single-leg differential onto $\Theta W$ is an isomorphism, and
the joint differential in (MCF14) has kernel the full $W$ diagonal.
Thus $K_{W,S;A}$ retains exactly the degree-zero diagonal specified
in (MCF36), including its zero-insertion and source-inclusion maps.
The original outer source,
every coefficient mask including the empty one, every present source
leg, and the exact arithmetic comparison kernel are simultaneously
specified by this diagram.

\subsection{Exact enlargement from offcritical support to all nontrivial packets}

The offcritical region $U$ throughout the preceding construction remains
unchanged.  A packet containing critical-line zeros has a larger original
source domain, which we now attach by its exact inclusions.  Put
\[
\begin{gathered}
\Omega=\{s\in\mathbb C:0<\Re s<1\},\qquad
L=\Omega\setminus U=\{s:\Re s=\tfrac12\},\\
T_XQ=\bigcup_{p\in\mathcal P_X}\ker(p(D):Q\to Q),\qquad
\mathscr B_X=q^{-1}(T_XQ),\qquad
C_X=[V\xrightarrow\Theta\mathscr B_X]\quad(X=U,\Omega,L),\\
\mathcal P_X=\{p\text{ monic}:\text{every root of }p\text{ belongs to }X\},
\qquad C_+=[V\xrightarrow\Theta\mathscr B].
\end{gathered}
\tag{MCF41}
\]
The polynomial $1$ is in every $\mathcal P_X$.  The notation $C_L$
denotes an algebraic source complex; no openness of the line $L$
or sheaf restriction to an open $L$ is required.  The same polynomial
range identities (MCF3) hold for every root in the open strip
$\Omega$, including its critical line.

The exact source inclusions, with their original degree-zero maps, are
\[
C_U[A]\xrightarrow{(1,\,\mathrm{inc})}C_\Omega[A]
       \xrightarrow{(1,\,\mathrm{inc})}C_+[A].
\tag{MCF41a}
\]
Indeed $\mathcal P_U\subseteq\mathcal P_\Omega$ gives
$T_UQ\subseteq T_\Omega Q$ and then
$\mathscr B_U\subseteq\mathscr B_\Omega\subseteq\mathscr B$.
All differentials are the same original $\Theta$ and the degree-zero
maps are the identity on $V^A$, proving both cochain identities.
Injectivity holds in each degree.  Cohomology in degree zero is zero
by injectivity of $\Theta$, and the induced degree-one sequence is
the literal sequence of submodules
$ (T_UQ)^A\hookrightarrow(T_\Omega Q)^A\hookrightarrow Q^A$.

Its full quotient has the exact arithmetic meaning
\[
T_\Omega Q=T_UQ\oplus T_LQ,\qquad
T_\Omega Q/T_UQ\xrightarrow{\sim}T_LQ.
\tag{MCF41b}
\]
Here is a direct proof without presuming a list of existing zeros.
For $u\in T_\Omega Q$, choose an annihilator $p$ in
$\mathcal P_\Omega$ and factor it as $p=p_Up_L$, retaining every
multiplicity, with roots of $p_U$ in $U$ and roots of $p_L$ in $L$.
Either factor can be $1$.  The factors have disjoint roots, so there
are polynomials $a,b$ with $ap_U+bp_L=1$.  Then
\[
u_U=b(D)p_L(D)u,\qquad u_L=a(D)p_U(D)u,
\qquad u=u_U+u_L,
\tag{MCF41c}
\]
and multiplication shows $p_U(D)u_U=p_L(D)u_L=0$.
If a vector is killed by a polynomial rooted in $U$ and by a
polynomial rooted in $L$, the same Bezout identity makes that vector
zero.  Thus the two submodules intersect in zero and the decomposition
is unique, independent of the chosen annihilator and Bezout
polynomials.  This proves (MCF41b), including its inverse: include
$T_LQ$ in $T_\Omega Q$ and take its class modulo $T_UQ$.

There is an exact original Mellin description of this quotient as well:
\[
T_XQ\xrightarrow{\sim}
\bigoplus_{\substack{\rho\in X\\g(\rho)=0}}
                     \mathcal O_\rho/(g),\qquad
[F]\longmapsto([\mathcal MF]_\rho)_\rho,
\quad X=U,L,\Omega.
\tag{MCF41d}
\]
To check the map, write $p(D)F=\Theta\phi$.  Then
$p(s)\mathcal MF(s)=g(s)H_\phi(s)$, so the Mellin class is zero
at every zero of $g$ outside the finite root set of $p$.  Replacing
$F$ by a theta boundary changes its Mellin transform by $gH$ and
does not change a local class.  For injectivity, suppose all local
classes at roots of $p$ vanish.  At a root with $g$ nonzero, $g$ is
a local unit; at a root of $g$, the vanishing class says that
$\mathcal MF$ is a multiple of $g$.  In either case write
$\mathcal MF=gB$ locally.  Cancellation in the local analytic domain
gives $H_\phi=pB$, so $j_pH_\phi=0$ at every required order.
By (MCF3), $\phi=p(D)\psi$ with $\psi\in V$.  Injectivity of $p(D)$
on $\mathscr B$ then yields $F=\Theta\psi$.

For surjectivity, take a finitely supported family of full local jets
$\alpha_\rho$ at actual zeros in $X$.  For its finite support $Z$ put
$h=\prod_{\rho\in Z}(s-\rho)^{m_\rho}$ with the complete actual
orders and choose the unique polynomial $P$ of degree less than
$\deg h$ with
$j_hP=j_h(g/h)^{-1}\alpha$.  The complete unit exists at each
specified centre and CRT supplies that polynomial.  The original
Euler range identity gives
\[
F_\alpha=S_h^{\mathscr B}\Theta(P(D)\phi_*),\qquad
\mathcal MF_\alpha=(g/h)P,\qquad h(D)F_\alpha\in\Theta V.
\tag{MCF41e}
\]
These equations give the specified full jets and zero classes at all
other zeros of $g$.  The empty family uses the stage $h=1$ and the
zero class.  Changing $P$ by $hR$ changes $F_\alpha$ by the original
boundary $\Theta(R(D)\phi_*)$.  This proves both inverse identities
in (MCF41d), retaining the entire $g/h$, every coordinate $s-\rho$,
and every multiplicity.  Under this isomorphism the decomposition
(MCF41b) is exactly zero insertion and restriction to the disjoint
offcritical and critical-line sets of full local factors.

Degreewise quotients of the actual complexes now give
\[
\begin{gathered}
0\longrightarrow C_U[A]\longrightarrow C_\Omega[A]
\longrightarrow (T_\Omega Q/T_UQ)^A[-1]\longrightarrow0,\\
0\longrightarrow C_\Omega[A]\longrightarrow C_+[A]
\longrightarrow (Q/T_\Omega Q)^A[-1]\longrightarrow0.
\end{gathered}
\tag{MCF41f}
\]
Their degree-zero quotients vanish because both inclusions are the
identity on $V^A$.  For the first degree-one quotient, the map
$\mathscr B_\Omega/\mathscr B_U\to T_\Omega Q/T_UQ$ induced by $q$
is onto, and has zero kernel because $qF\in T_UQ$ is exactly
$F\in\mathscr B_U$.  The second quotient is proved by the identical
calculation with $\mathscr B/\mathscr B_\Omega$ and $Q/T_\Omega Q$.
These prove the displayed short exact sequences on the original
source domains.  In particular a nonzero critical local factor
survives in the first quotient; it has not been discarded by working
with the original offcritical subcomplex.

For every actual nontrivial full packet $h$ in $\Omega$, including a
critical pair or a mixture of critical and offcritical centres,
the formulas (MCF4)--(MCF5) and (MCF22)--(MCF27) have their unchanged
values and their source now lies in $\mathscr B_\Omega$ because
$h(D)s_hu\in\Theta V$.  Their proofs used (MCF3) in the open strip
and the full local unit; they did not use $\Re\rho\ne1/2$.
Thus they apply to that precise larger domain.  Their restriction
to packets rooted in $U$ agrees literally with the earlier formulas
through (MCF41a).  Every packet system includes $h=1$ with $E_1=0$;
no nonempty offcritical packet is assumed to make the enlarged
construction or its embeddings.

Finally retain every coefficient and outer label in this enlargement.
Replacing only the central term of (MCF34) by $\mathscr B_X^A$
defines $C^X_{W,S;A}$ and $Q_{W,X}=\mathscr B_X/\Theta W$.
The maps from $X=U$ to $X=\Omega$ are the identity on the original
source legs and the inclusion on the central term, and their
degree-one quotient is again $(T_\Omega Q/T_UQ)^A$:
\[
0\longrightarrow Q_{W,U}^A\longrightarrow Q_{W,\Omega}^A
  \longrightarrow (T_\Omega Q/T_UQ)^A\longrightarrow0.
\tag{MCF41g}
\]
Indeed $\Theta W\subseteq\mathscr B_U$; quotienting the two nested
central subspaces by that same subspace proves injectivity and the
stated quotient.  The proper-source kernel $(V/W)^A$ of (MCF36)
remains attached on each side.  Empty coefficient support still
retains $(W,S)$ as in (MCF35).  The enlarged maps commute with zero
insertion, the specified slot-local comparisons, dilation and its
proper-source transport, and coefficient Frobenius, since their
source maps are identities and their central maps are the original
coordinatewise inclusions.  This attaches all actual nontrivial
packets while preserving the original offcritical construction and
its full comparison quotient.

\subsection{Literal comparison with the earlier finite source and its cocycle}

The earlier source (TC1) writes the spectral polynomial coordinate as
$t$ and the same residue representative as $R_Z(u)$; the coordinate
isomorphism sends $t$ to $s$, with $h_Z(t)$ sent to $h(s)$, every
power $(t-\rho)^r$ sent to $(s-\rho)^r$, and the full unit unchanged.
Consequently $R_Z(u)$ is $P_h(u)$ under this specified isomorphism.
For the actual source functions, not only their quotient classes,
\[
R_{\mathrm{ref}}=s_h:E_h\longrightarrow\mathscr B_X,
\qquad X=U\text{ or }\Omega\text{ as the packet requires}.
\tag{MCF42}
\]
Indeed applying $h(D)$ to either representative gives
$\Theta(P_h(u)(D)\phi_*)$.  The difference lies in $\mathscr B$
and is killed by the injective $h(D)$, hence vanishes.  This proves
the asserted equality with the original units and domains.

Put $t=\log a$, retaining its orientation when $a<1$, and define
\[
c_{t,h}(u)=\int_0^t\mathcal R_{e^{t-r}}\phi_*\,
                  \ell_h(e^{rA_h}u)\,dr:E_h\longrightarrow V.
\tag{MCF43}
\]
The integral converges in the original Schwartz topology.  In fact
on the compact interval joining $0$ and $t$, every derivative of
$\mathcal R_{e^{t-r}}\phi_*$ is a continuously varying rescaled
derivative of the same Schwartz function, with the scaling parameter
in a compact subset of $(0,\infty)$.  For every pair $N,k$, its
seminorm $\sup_x(1+|x|)^N|\partial_x^k(\cdot)|$ is uniformly
bounded there.  The finite matrix $e^{rA_h}$ is smooth and bounded
there too.  Riemann sums therefore converge in every seminorm in
the complete Schwartz space; parity and the two zero moments pass
to the limit by continuity, so the limit is in $V$.

Differentiate the actual $\mathscr B$-valued function
$\mathcal R_{e^{t-r}}s_he^{rA_h}u$.  Its derivative is
$-\mathcal R_{e^{t-r}}(Ds_h-s_hA_h)e^{rA_h}u$.
The same seminorm argument on $\mathscr B$, followed by (MCF5),
allows the fundamental theorem of calculus.  It gives
\[
\Theta c_{\log a,h}=\mathcal R_as_h-s_hA_{a,h},\qquad
\Psi_{a,h}=c_{\log a,h}.
\tag{MCF44}
\]
The second equality follows by comparing with (MCF24) and canceling
the injective map $\Theta:V\to\mathscr B$.  Thus the previous
oriented integral and the Euler-inverse primitive are the same map.
Their packet and coefficient-face naturality are literal, by
(MCF27) and coordinatewise integration.  The joint primitive is
$(c/2,-\widehat c/2)$; the difference between the two single-leg
primitives $(c,0)$ and $(0,-\widehat c)$ is the full cycle
$(c,\widehat c)$.  No diagonal is discarded by this comparison.

The earlier dual two-leg complex is obtained by applying algebraic
$\operatorname{Hom}_{\mathbb C}(-,\mathbb C)$, so the map induced
by $C_U\to C_+$ is restriction $\mathscr B^\vee\to\mathscr B_U^\vee$
in degree $-1$ and identity on both $V^\vee$ source coordinates in
degree zero.  Its differential, in the same order of legs, is
\[
\ell\longmapsto(\ell\Theta,-\ell\Theta\mathcal F),
\qquad \mathcal F\phi=\widehat\phi.
\tag{MCF45}
\]
Because $\Theta V\subseteq\mathscr B_U$, restriction followed by
this differential equals the differential followed by the identity.
Thus this is an actual cochain map.  It induces restriction
$Q^\vee\to(T_UQ)^\vee$ on degree $-1$ cohomology.  At a joint
source the two diagonal parameters $v\mapsto(v,\mathcal Fv)$ and
$w\mapsto(\mathcal Fw,w)$ agree under $w=\mathcal Fv$. Define
$\mathcal F^\vee:V^\vee\to V^\vee$ by
$\mathcal F^\vee(\ell)=\ell\circ\mathcal F$; the maps
$\pi_v,\pi_w:V^\vee\oplus V^\vee\to V^\vee$ are
\[
\pi_v(\alpha,\beta)=\alpha+\beta\mathcal F,
\quad \pi_w(\alpha,\beta)=\alpha\mathcal F+\beta,
\quad \pi_v=\mathcal F^\vee\circ\pi_w.
\tag{MCF46}
\]
These identities follow by evaluation on the displayed diagonals
and $\mathcal F^2=1$.  The original actions intertwine because
$\mathcal F\mathcal R_a=a\mathcal R_{1/a}\mathcal F$.
For a proper source $W$, dualizing $W\hookrightarrow V$ gives
restriction of diagonal functionals to $W$.  For coefficient faces,
dualizing zero insertion gives coordinate restriction.  These are
algebraic dual statements; no continuous splitting for arbitrary
nonclosed $W$ is asserted.
\subsection{The shifted coefficient quotient on the original theta complexes}

Fix one of the actual full packets $h$ in $\Omega$ and a member
$\rho$ with $m=m_\rho\geq2$.  The calculation is indexed by precisely
those members; it requires no assertion that the indexing set is
nonempty.  The coordinate isomorphism from the $x$ presentation
in (MW23) to this $s$ presentation sends $x$ to $s$ and
$e_\rho(x-\rho)^r$ to $e_\rho(s-\rho)^r$, with coefficient one
and inverse $s\mapsto x$.  It carries the full defining
polynomial, its CRT idempotents, and every retained jet order.
Use the original local subspace $E_\rho=e_\rho E_h$,
$J_\rho=(s-\rho)^2E_\rho$, and $C_\rho=E_\rho/J_\rho$ of
(MW23).  Write $\sigma_\rho$ for the restriction of the original
injective $\sigma_h:E_h\to Q$.  Define actual subspaces of the
original central source by
\[
\mathscr B_\rho=q^{-1}(\sigma_\rho E_\rho),\qquad
\mathscr B_{J_\rho}=q^{-1}(\sigma_\rho J_\rho),\qquad
\Theta V\subseteq\mathscr B_{J_\rho}\subseteq\mathscr B_\rho
       \subseteq\mathscr B_\Omega .
\tag{MCF47}
\]
The last inclusion follows because every displayed quotient class
lies in the $h(D)$-torsion, and $h$ is rooted in $\Omega$.
The literal direct-sum equalities are
\[
\mathscr B_\rho=\Theta V\oplus s_h E_\rho,\qquad
\mathscr B_{J_\rho}=\Theta V\oplus s_h J_\rho.
\tag{MCF48}
\]
For the first, if $qF=\sigma_\rho x$, then $F-s_hx\in\Theta V$;
injectivity of $\sigma_\rho$ makes this $x$ unique.  If
$\Theta v+s_hx=0$, quotienting gives $x=0$, then injectivity of
$\Theta$ gives $v=0$.  The same argument with $x\in J_\rho$
proves the second equality.  Every representative here is the
original $s_hx=P_h(x)(D)F_h$ with the full unit
$j_h(g/h)^{-1}$, including its components at all selected centres.

There is an actual short exact sequence of source complexes
\[
0\longrightarrow[V\xrightarrow{\Theta}\mathscr B_{J_\rho}]
\xrightarrow{(1,\mathrm{inc})}
[V\xrightarrow{\Theta}\mathscr B_\rho]
\longrightarrow C_\rho[-1]\longrightarrow0.
\tag{MCF49}
\]
The final map is zero in degree zero.  In degree one it sends
$F$ to $[x]$, where the unique $x\in E_\rho$ satisfies
$qF=\sigma_\rho x$.  It is linear, vanishes on $\Theta V$,
is onto because $s_hx$ is a representative of every $[x]$,
and has kernel exactly $\mathscr B_{J_\rho}$.  These assertions
prove both the cochain identity and degreewise exactness.
Its inverse on the degree-one quotient sends $[x]$ to
$[s_hx]\in\mathscr B_\rho/\mathscr B_{J_\rho}$; replacing $x$
by $x+j$, $j\in J_\rho$, changes its representative by the
actual element $s_hj\in\mathscr B_{J_\rho}$.  Both inverse
identities follow from (MCF48).  Thus the length-two quotient
used in (MW24)--(MW27) and (MRE34)--(MRE42) has this explicit
original-source realization.

The arithmetic actions in this exact sequence also have their
original types.  Both $E_\rho$ and $J_\rho$ are invariant under
$A_h$, $N_\rho=A_h-\rho I$ on $E_\rho$, and each $A_{a,h}$.
The full-source intertwining (MCF5), (MCF21) makes
$D$, $D-\rho$, and $\mathcal R_a$ preserve the corresponding
spaces in (MCF47).  Their maps on degree zero are respectively
$D$, $D-\rho$, and $\mathcal R_a$ on $V$, and their maps on
degree one use the same original operators on functions.
Since each commutes with $\Theta$, they are cochain maps.
Under the final map in (MCF49) they give exactly the quotient
$\overline A_h$, $\overline N_\rho$, and $\overline A_{a,h}$
of (MW25), with every surviving coefficient retained.  The
full higher powers remain in the original source operators and
in the subcomplex on the left.

The induced filtration is realized before taking cohomology.
For every integer $j$, set
\[
\begin{gathered}
\mathscr B_{\rho,j}=q^{-1}(\sigma_\rho(M_jE_\rho)),\qquad
\mathscr B_{J_\rho,j}
  =q^{-1}(\sigma_\rho(J_\rho\cap M_jE_\rho))
  =\mathscr B_{J_\rho}\cap\mathscr B_{\rho,j},\\
0\longrightarrow[V\xrightarrow{\Theta}\mathscr B_{J_\rho,j}]
\longrightarrow[V\xrightarrow{\Theta}\mathscr B_{\rho,j}]
\longrightarrow\overline M_jC_\rho[-1]\longrightarrow0.
\end{gathered}
\tag{MCF50}
\]
The intersection equality uses injectivity of $\sigma_\rho$.
For the exact sequence, the last map of (MCF49) sends
$\mathscr B_{\rho,j}$ onto $\pi_\rho(M_jE_\rho)$, because the
representatives $s_hx$ exist for every $x\in M_jE_\rho$.
Its kernel is precisely the displayed intersection.  The
degree-zero map remains identity and the last differential is
zero, proving every case.  Consequently the quotient filtration
is exactly (MW24), with degrees $m-1$ and $m-3$, hence with its
retained shift $m-2$ relative to the displayed length-two
centered filtration.

For all sufficiently negative $j$, the two central terms in
(MCF50) equal $\Theta V$.  Their source complexes then remain
$[V\xrightarrow{\Theta}\Theta V]$, with its actual invertible
differential; no assertion that these source complexes are
literally zero is made.  They have zero cohomology because
$\Theta$ is injective and onto its displayed image.  At every
step the original $D$ preserves $\mathscr B_{\rho,j}$, and
$D-\rho$ sends it into $\mathscr B_{\rho,j-2}$ since its
quotient action is precisely $N_\rho$.  This proves the
filtered source-map types as well as the induced filtered
cohomology identities.

For each proper original $W\subseteq V$, replace the degree-zero
term in (MCF49)--(MCF50) by the outer source
$\bigoplus_{+\in S}W\oplus\bigoplus_{-\in S}\widehat W$ with
the original signed theta differential.  Use that identical
source in the left and middle complexes.  The same proof of
degreewise exactness applies, and in degree one their individual
cohomology groups have the exact sequences
\[
\begin{gathered}
0\longrightarrow V/W\longrightarrow
\mathscr B_{\rho,j}/\Theta W\longrightarrow M_jE_\rho
\longrightarrow0,\\
0\longrightarrow V/W\longrightarrow
\mathscr B_{J_\rho,j}/\Theta W\longrightarrow
J_\rho\cap M_jE_\rho\longrightarrow0.
\end{gathered}
\tag{MCF51}
\]
In both lines $[v]\mapsto[\Theta v]_W$ is the original kernel
map, and the last map is $\sigma_\rho^{-1}q$ on its indicated
image.  Injectivity, surjectivity, and the kernel computation
are the same as in (D11).  Thus the low-filtration degree-one
cohomology is $V/W$ at a proper source, while the joint
degree-zero diagonal is $W$; both remain present.
For $D$ use $DW\subseteq W'$ and the specified source codomain
$W'$ as in (D14).  For $D-\rho$ use $(D-\rho)W\subseteq W'$;
its minus source operator is $1-D-\rho$, because
$\mathcal F(D-\rho)=(1-D-\rho)\mathcal F$.
The choice $W'=W+DW$ admits both maps.  Dilation uses
$\mathcal R_aW$.  Apply the same proof independently at every
coefficient mask, including its retained empty outer label.
Zero insertion and coefficient Frobenius commute with the
displayed quotient maps by linearity and their original signed
coordinate permutation.  This transports the strict shifted
coefficient quotient to the original mixed source system
without removing its actual comparison kernels.
